\chapter{Introduzione ai sistemi informatici}
    \section{introduzione}
    \section{l'informazione e la sua codifica}
    \section{l'elaborazione e la struttura dell'informazione}
    \section{i linguaggi per la formalizzazione dell'informazione}
    \section{la trasmissione dell'informazione}
    \section{le infrastrutture hardware}
        l'architettura di riferimento
        l'esecutore
            la struttura del processore
            le CPU pipeline
            evoluzione delle CPU
        la memoria
            parametri per la descrizione di una memoria
            organizzazione gerarchica della memoria
        i dispositivi per le memorie di massa
            i dischi e i nastri magnetici
            i dischi ottici
        l'interfaccia ingresso/uscita
            le operazioni I/O a controllo del programma
            le operazioni I/O gestite tramite interruzione
            le operazioni I/O con accesso diretto alla memoria
        le principali periferiche
            il video e la tastiera
            i dispositivi di puntamento
            le stampanti
            le periferiche per la gestione di immagini, suoni e movimenti
    \section{le infrastrutture software}
        le funzioni del sistema operativo
            gli elementi di un sistema operativo
            l'evoluzione dei sistemi operativi
            i modelli organizzativi dei sistemi operativi
        i processi
            la gestione dei processi
            il modello client-server
            i modelli ibridi
        la gestione della memoria
            paginazione e memoria virtuale
            la segmentazione della memoria
        la gestione delle periferiche
        il file system
            la localizzazione dei dati
            i servizi di base
        la gestione della rete
    \section{le infrastrutture di rete}
        i mezzi di trasmissione
        la tecnologia di trasmissione
        le reti geografiche
            la comunicazione a circuito
            la comunicazione a pacchetto
            ISDN (Integrated Services Digital Network)
            DSL (Digital Subscriber Line - Asymmetric, Very-high-bit-rate, V2)
            WiMAX
        le reti locali
            i mezzi di trasmissione e le topologie di rete
            i metodi di accesso
            le tipologie di rete
            l'interconnessione di LAN
        l'architettura del software di rete
            l'architettura a livelli di un insieme di protocolli
            i principi di funzionamento di TCP/IP
            indirizzi numerici e indirizzi simbolici
    \section{le applicazioni}
        l'interazione con l'utente 
        applicazioni come sistemi 
        le applicazioni come sistemi distribuiti 
            le applicazioni web 
        le applicazioni per dispositivi mobili 
        le tecnologie per il World Wide Web
            il WWW e l'identificazione delle risorse
            l'utilizzo di un'applicazione web
            le tecnologie per le applicazioni web-oriented
        il problema della sicurezza
% ----------------------------------------------------
\chapter{Reti di calcolatori e internet}
    \section{reti di calcolatori e internet}
    \section{livello applicazione}
    \section{livello trasporto}
    \section{livello rete}
    \section{livello collegamento: collegamenti, reti di accesso e locali}
    \section{reti multimediali}
        applicazioni multimediali di rete
            proprietà del video
            proprietà dell'audio
            tipi di applicazioni multimediali
        streaming di video registrato
            streaming UDP
            streaming HTTP
            streaming adattivo e DASH
            reti per la distribuzione di contenuti (CDN)
        Voice-over-IP
            limiti del servizio best-effort IP
            rimozione del jitter al ricevente
            recupero dei pacchetti persi
            un caso di studio: VoIP con Skype
        protocolli per applicazioni in tempo reale
            RTP (Realtime Transport Protocol)
            SIP (Session Initiation Protocol)
        supporto di internet alle applicazioni multimediali
            dimensionamento delle reti best-effort
            fornitura di più classi di servizio
            diffserv (differentiated services)
            fornire garanzie di qualità del servizio (QoS) per ogni connessione: prenotazione delle risorse e ammissione delle chiamate
    \section{sicurezza nelle reti}
        sicurezza di rete
        principi di crittografia
            crittografia a chiave simmetrica
            crittografia a chiave pubblica
        integrità dei messaggi e firma digitale
            funzioni hash crittografiche
            codice di autenticazione dei messaggi
            firme digitali
        autenticazione
            protocollo auth ap1.0
            protocollo auth ap2.0
            protocollo auth ap3.0
            protocollo auth ap4.0
        rendere sicura la posta elettronica
            email secure
            PGP (pretty good privacy; asym)
        rendere sicure le connessioni TCP:SSL
            quadro generale
            un quadro più completo
        sicurezza a livello di rete: IPsec
            IPsec e le reti virtuali private VPN
            i protocolli AH (auth header) e ESP (encapsulated sec. protc)
            associazioni di sicurezza
            il datagramma IPsec
            IKE: gestione delle chiavi in IPsec (internet key exchange)
        sicurezza nelle LAN wireless
            wired equivalent privacy (WEP)
            IEEE 802.11i --> WiFi Protected Access 2 (WPA2 / WPA3)
        sicurezza operativa: firewall e sistemi di rilevamento delle intrusioni
            firewall
            sistemi di rilevamento delle intrusioni
    \section{gestione della rete}
        che cosa vuol dire gestire una rete?
        infrastruttura di gestione standard di internet
            struttura delle informazioni di gestione
            base di dati gestionali: MIB (management information base)
            operazioni del protocollo SNMP (simple network management protocol)
            sicurezza e amministrazione
        ASN.1 (Abstract Syntax Notation One)
% ----------------------------------------------------
\chapter{Reti di calcolatori (Sesta Edizione)}
    \section{Sicurezza delle reti}
        fondamenti di sicurezza delle reti
            principi fondamentali della sicurezza
            principi fondamentali degli attacchi
            dalle minacce alla sicurezza alle soluzioni
        elementi fondamentali di un attacco
            reconnaissance
            sniffing e snooping (con un pizzico di spoofing)
            spoofing (oltre ARP)
            disruption (interruzione)
        firewall e sistemi di rilevamento delle intrusioni
            firewall
            rilevamento e prevenzione delle intrusioni
        crittografia
            introduzione alla crittografia
            due principi crittografici fondamentali
            cifrari a sostituzione
            cifrari a trasposizione
            blocchi monouso
        algoritmi a chiave simmetrica
            DES - Data Encryption Standard
            AES - Advanced Encryption Standard
            modalità di cifratura
        algoritmi a chiave pubblica
            RSA
            altri algoritmi a chiave pubblica
        firme digitali
            firme a chiave simmetrica
            firme a chiave pubblica
            message digest
            attacco del compleanno
        gestione delle chiavi pubbliche
            certificati
            X.509
            infrastrutture a chiave pubblica
        protocolli di autenticazione
            autenticazione basata su un segreto condiviso
            come stabilire una chiave condivisa: lo scambio di chiave Diffie-Hellman
            autenticazione usando un centro di distribuzione delle chiavi
            autenticazione con kerberos
            autorizzazione con crittografia a chiave pubblica
        sicurezza delle comunicazioni
            IPsec
            VPN - Reti Virtuali Private
            Sicurezza Wireless
        sicurezza della posta elettronica
            PGP - Pretty Good Privacy
            S/MIME (Secure/Multipurpose Internet Mail Extensions)
        sicurezza del web
            minacce alla sicurezza
            sicurezza nell'attribuzione dei nomi e DNSSEC
            sicurezza del livello trasporto
            sicurezza con codici non affidabili
        aspetti sociali
            privacy
            libertà di parola
            copyright
% ----------------------------------------------------
\chapter{The Nature of Code}
    \section{Neural Networks}
        introducing Artifical Neural Networks
        the Perceptron
        putting the "network" in neural network
        machine learning with ml5.js
        building a gesture classifier
    \section{Neuroevolution}
        reinforcement learning
        evolving neural network is neat!
        coding flappy bird
        neuroevolutionary flappy bird
        steering the neuroevolutionary way
        a neuroevolutionary ecosystem
        